\documentclass[aspectratio=169, 10pt]{beamer}

\usetheme{Madrid}
\usefonttheme{professionalfonts}

\usepackage[utf8]{inputenc}
\usepackage[indonesian]{babel}
\usepackage{amsmath,amsfonts,amssymb}
\usepackage{graphicx}
\usepackage{booktabs}
\usepackage{colortbl}
\usepackage{listings}
\usepackage{xcolor}

\definecolor{mainblue}{RGB}{24, 43, 73}
\definecolor{accentblue}{RGB}{41, 128, 185}
\definecolor{accentgold}{RGB}{230, 126, 34}
\definecolor{darkgray}{RGB}{44, 62, 80}
\definecolor{lightbg}{RGB}{248, 250, 252}
\definecolor{bordergray}{RGB}{200, 210, 220}
\definecolor{codebg}{RGB}{240, 243, 246}
\definecolor{codegreen}{RGB}{39, 174, 96}
\definecolor{tableheader}{RGB}{24, 43, 73}

\setbeamercolor{structure}{fg=mainblue}
\setbeamercolor{palette primary}{bg=mainblue, fg=white}
\setbeamercolor{palette secondary}{bg=accentblue, fg=white}
\setbeamercolor{palette tertiary}{bg=mainblue!80, fg=white}
\setbeamercolor{title}{fg=white, bg=mainblue}
\setbeamercolor{frametitle}{fg=white, bg=mainblue}
\setbeamercolor{block title}{bg=mainblue, fg=white}
\setbeamercolor{block body}{bg=lightbg, fg=darkgray}
\setbeamercolor{block title alerted}{bg=accentgold, fg=white}
\setbeamercolor{block body alerted}{bg=accentgold!10, fg=darkgray}

\lstset{
    language=Python,
    backgroundcolor=\color{codebg},
    basicstyle=\ttfamily\scriptsize\color{darkgray},
    keywordstyle=\bfseries\color{accentblue},
    stringstyle=\color{accentgold},
    commentstyle=\itshape\color{codegreen},
    showstringspaces=false,
    breaklines=true,
    frame=single,
    rulecolor=\color{bordergray},
    numbers=none
}

\title[\textbf{Deret Waktu (Minggu 7)}]{\textbf{Discrete Kalman Filter untuk Estimasi dan Filtering}}
\subtitle{Representasi Ruang Keadaan (State-Space), Denoising Sensor, dan Estimasi SoC Baterai}
\author[\textbf{Ir. Novalio Daratha, Ph.D.}]{\textbf{Ir. Novalio Daratha, S.T., M.Sc., Ph.D.}}
\institute[FT UNIB]{Program Studi Teknik Elektro -- Fakultas Teknik\\Universitas Bengkulu}
\date{Minggu 7 / Semester Ganjil 2026}

\begin{document}

\begin{frame}[plain]
    \titlepage
\end{frame}

\begin{frame}{Agenda Perkuliahan Minggu ke-7}
    \begin{columns}[T]
        \begin{column}{0.48\textwidth}
            \begin{block}{Capaian Pembelajaran (Sub-CPMK 7)}
                Mahasiswa mampu memformulasikan model sistem dinamis ke dalam ruang keadaan (\textit{State-Space}), menurunkan algoritma Discrete Kalman Filter, serta mengimplementasikannya untuk estimasi parameter dan penyaringan derau sensor elektro.
            \end{block}
        \end{column}
        \begin{column}{0.48\textwidth}
            \begin{enumerate}
                \item \textbf{Representasi Ruang Keadaan (State-Space)}
                \item \textbf{Derau Proses ($Q$) vs Derau Pengukuran ($R$)}
                \item \textbf{Algoritma 2-Tahap: Prediksi dan Koreksi}
                \item \textbf{Kalman Gain ($K_k$) dan Konvergensi Kovarians}
                \item \textbf{Studi Kasus: Estimasi State-of-Charge (SoC)}
                \item \textbf{Review Komprehensif Materi UTS}
            \end{enumerate}
        \end{column}
    \end{columns}
\end{frame}

\begin{frame}{Representasi Ruang Keadaan Diskrit}
    Sistem dinamis linier stokastik dimodelkan dalam dua persamaan:
    \begin{block}{1. Persamaan State (Model Fisik Sistem)}
        \begin{equation}
            \mathbf{x}_k = \mathbf{A}_k \mathbf{x}_{k-1} + \mathbf{B}_k \mathbf{u}_{k-1} + \mathbf{w}_{k-1}, \quad \mathbf{w}_k \sim \mathcal{N}(0, \mathbf{Q}_k)
        \end{equation}
    \end{block}
    \begin{block}{2. Persamaan Pengukuran (Sensor)}
        \begin{equation}
            \mathbf{z}_k = \mathbf{H}_k \mathbf{x}_k + \mathbf{v}_k, \quad \mathbf{v}_k \sim \mathcal{N}(0, \mathbf{R}_k)
        \end{equation}
    \end{block}
    \begin{columns}[T]
        \begin{column}{0.48\textwidth}
            $\mathbf{x}_k$: Vektor status tersembunyi (misal SoC baterai, fluks motor).\\
            $\mathbf{z}_k$: Sinyal pengukuran nyata dari sensor yang bising.
        \end{column}
        \begin{column}{0.48\textwidth}
            $\mathbf{Q}$: Matriks kovariansi ketidakpastian model.\\
            $\mathbf{R}$: Matriks kovariansi derau instrumen sensor.
        \end{column}
    \end{columns}
\end{frame}

\begin{frame}{Algoritma Siklus 2-Tahap Discrete Kalman Filter}
    \begin{columns}[T]
        \begin{column}{0.48\textwidth}
            \begin{block}{Tahap 1: Prediksi (Time Update)}
                Memproyeksikan status ke depan:
                \begin{equation*}
                \begin{aligned}
                    \hat{\mathbf{x}}_k^- &= \mathbf{A}\hat{\mathbf{x}}_{k-1} + \mathbf{B}\mathbf{u}_{k-1} \\
                    \mathbf{P}_k^- &= \mathbf{A}\mathbf{P}_{k-1}\mathbf{A}^T + \mathbf{Q}
                \end{aligned}
                \end{equation*}
                $\hat{\mathbf{x}}_k^-$: Estimasi a priori.\\
                $\mathbf{P}_k^-$: Kovariansi error prediksi.
            \end{block}
        \end{column}
        \begin{column}{0.48\textwidth}
            \begin{block}{Tahap 2: Koreksi (Measurement Update)}
                Memperbaiki estimasi dengan data sensor:
                \begin{equation*}
                \begin{aligned}
                    \mathbf{K}_k &= \mathbf{P}_k^-\mathbf{H}^T(\mathbf{H}\mathbf{P}_k^-\mathbf{H}^T + \mathbf{R})^{-1} \\
                    \hat{\mathbf{x}}_k &= \hat{\mathbf{x}}_k^- + \mathbf{K}_k(\mathbf{z}_k - \mathbf{H}\hat{\mathbf{x}}_k^-) \\
                    \mathbf{P}_k &= (\mathbf{I} - \mathbf{K}_k\mathbf{H})\mathbf{P}_k^-
                \end{aligned}
                \end{equation*}
                $\mathbf{K}_k$: Kalman Gain penentu bobot optimal.
            \end{block}
        \end{column}
    \end{columns}
\end{frame}

\begin{frame}{Perilaku Dinamis Kalman Gain}
    Bagaimana Kalman Filter menyeimbangkan antara model teori vs data sensor?
    \begin{columns}[T]
        \begin{column}{0.48\textwidth}
            \begin{block}{Kondisi 1: Sensor Sangat Akurat ($\mathbf{R} \to 0$)}
                \begin{equation}
                    \lim_{\mathbf{R} \to 0} \mathbf{K}_k = \mathbf{H}^{-1} \quad \Longrightarrow \quad \hat{\mathbf{x}}_k \approx \mathbf{H}^{-1}\mathbf{z}_k
                \end{equation}
                Filter mempercayai pembacaan sensor secara penuh dan mengabaikan model teoritis.
            \end{block}
        \end{column}
        \begin{column}{0.48\textwidth}
            \begin{block}{Kondisi 2: Model Sempurna ($\mathbf{R} \to \infty$)}
                \begin{equation}
                    \lim_{\mathbf{R} \to \infty} \mathbf{K}_k = 0 \quad \Longrightarrow \quad \hat{\mathbf{x}}_k \approx \hat{\mathbf{x}}_k^-
                \end{equation}
                Filter mengabaikan data sensor yang sangat bising dan mengandalkan perambatan model fisik.
            \end{block}
        \end{column}
    \end{columns}
    \vspace{0.2cm}
    \begin{center}
        Kalman Filter adalah Best Linear Unbiased Estimator (BLUE)!
    \end{center}
\end{frame}

\begin{frame}[fragile]{Implementasi Python: Kalman Filter 1-Dimensi}
\begin{lstlisting}
import numpy as np
import matplotlib.pyplot as plt

# Parameter simulasi sensor tegangan
A, H = 1.0, 1.0
Q = 1e-5  # Derau proses rendah (tegangan sistem stabil)
R = 0.04  # Derau sensor tinggi (standar deviasi ~ 0.2 V)

# Inisialisasi
x_hat, P = 0.0, 1.0
estimates, raw_noisy = [], df['voltage_noisy'].values

for z in raw_noisy:
    # 1. Time Update
    x_hat_minus = A * x_hat
    P_minus = A * P * A + Q
    
    # 2. Measurement Update
    K = P_minus * H / (H * P_minus * H + R)
    x_hat = x_hat_minus + K * (z - H * x_hat_minus)
    P = (1 - K * H) * P_minus
    estimates.append(x_hat)

plt.plot(raw_noisy, label='Sensor Bising (z)', alpha=0.5)
plt.plot(estimates, label='Estimasi Kalman (x_hat)', color='red', linewidth=2)
plt.legend(); plt.show()
\end{lstlisting}
\end{frame}

\begin{frame}{Rangkuman Paruh Pertama Semester (Persiapan UTS)}
    \begin{table}[h!]
    \centering
    \scriptsize
    \renewcommand{\arraystretch}{1.1}
    \begin{tabular}{|c|l|l|}
    \hline
    \rowcolor{tableheader}
    \color{white}\textbf{Minggu} & \color{white}\textbf{Fokus Materi Pokok} & \color{white}\textbf{Aplikasi Rekayasa Sistem Tenaga} \\ \hline
    1 & Dekomposisi Deret Waktu & Pemisahan tren dan variasi musiman beban harian PLN \\ \hline
    2 & Uji Stasioneritas & Uji ADF/KPSS dan eliminasi unit root pada daya aktif \\ \hline
    3 & Pemulusan Eksponensial & Peramalan cepat beban jangka pendek (Holt-Winters) \\ \hline
    4 & Pemodelan ARMA & Identifikasi sistem linier dan kestabilan bidang Z \\ \hline
    5 & Metodologi Box-Jenkins & Model musiman SARIMA peramalan 168 jam ke depan \\ \hline
    6 & Model Eksogen SARIMAX & Mengintegrasikan ramalan suhu/cuaca pada peramalan beban \\ \hline
    7 & Discrete Kalman Filter & Penyaringan derau sensor telemetri dan estimasi SoC baterai \\ \hline
    \rowcolor{lightbg}
    \textbf{8} & \multicolumn{2}{l|}{\textbf{UJIAN TENGAH SEMESTER (UTS) -- Evaluasi Kognitif CPMK-1 dan CPMK-2}} \\ \hline
    \end{tabular}
    \end{table}
\end{frame}

\begin{frame}[plain]
    \begin{center}
        \vspace{1.5cm}
        {\Huge \color{mainblue} \textbf{Terima Kasih dan Sukses UTS}}\\[0.4cm]
        {\large Selamat Belajar untuk Ujian Tengah Semester Minggu Depan!}\\[1.2cm]
        \textbf{Laboratorium Sistem Tenaga dan Komputasi Cerdas}\\
        Jurusan Teknik Elektro -- Universitas Bengkulu
    \end{center}
\end{frame}

\end{document}
